\documentclass{article}

\usepackage[T1]{fontenc}
\usepackage{multicol}
\usepackage{calc}
\usepackage{ifthen}
\usepackage{epsfig}
%\usepackage{bookman}
%\usepackage{booktabs}
\usepackage{fontspec}
\usepackage{upquote} % Uncurly the quotes
\usepackage{etoolbox} % for backquote fix
\usepackage{mdframed} % nice frames
\usepackage[nobottomtitles]{titlesec} % better titles
\usepackage{enumitem}
\usepackage{minted2}
\usepackage[a4paper,landscape,margin=1cm,includehead,headsep=2ex,nofoot]{geometry}
\usepackage{fancyhdr}
\usepackage{array}
\usepackage[document]{ragged2e}

\usemintedstyle{tango}

\hyphenpenalty=10000
\exhyphenpenalty=10000

\setsansfont{DejaVu Sans}[Scale=0.9]
\setmonofont{DejaVu Sans Mono}[Scale=0.9]

\setlist[itemize]{noitemsep, topsep=0pt}

\renewcommand{\MintedPygmentize}{../../_build/default/compiler/catala.exe pygmentize}

% fancyvrb and mdframed stuff
\newmdenv[bottomline=false,rightline=false]{topleft}
\newmdenv[topline=false,leftline=false]{bottomright}
\DefineVerbatimEnvironment{Verbacorner}{Verbatim}{}
\surroundwithmdframed[bottomline=false,rightline=false]{Verbacorner}
\fvset{commandchars=@\{\}}

\ifthenelse{\lengthtest { \paperwidth = 11in}}
	{ \geometry{top=.5in,left=.5in,right=.5in,bottom=.5in} }
	{\ifthenelse{ \lengthtest{ \paperwidth = 297mm}}
		{\geometry{top=1cm,left=1cm,right=1cm,bottom=1cm} }
		{\geometry{top=1cm,left=1cm,right=1cm,bottom=1cm} }
	}
\pagestyle{empty}

% Don't print section numbers
\setcounter{secnumdepth}{0}

\setlength{\parindent}{0pt}
\setlength{\parskip}{1ex}

\pagestyle{fancy}
\newcommand{\makeheader}[3]{
  \fancyhead[L]{
    \begin{minipage}{30pt}
      \vspace{-3pt}
     \includegraphics[width=\textwidth]{logo.pdf}
    \end{minipage}
  }
  \fancyhead[C]{
   \LARGE{#1}
  }
  \fancyhead[R]{
   {\scriptsize
    #2 $\cdot$ Révision \#1 $\cdot$ \copyright\ \the\year$\quad$
    %% \begin{minipage}{30pt}
    %%  \vspace{-3pt}
    %%  \includegraphics[width=\textwidth]{logo.pdf} %defined at bottom
    %% \end{minipage}
   }}
  \fancyfoot[L]{}
  \fancyfoot[C]{}
  \fancyfoot[R]{}
}
\renewcommand\headrule{}

\newenvironment{absolutelynopagebreak}
  {\par\nobreak\vfil\penalty0\vfilneg
   \vtop\bgroup}
  {\par\xdef\tpd{\the\prevdepth}\egroup
    \prevdepth=\tpd}

\newcommand{\vb}{\usefont{T1}{DejaVuSansMono-TLF}{b}{n}}

\newenvironment{libcomments}{\begin{itemize}[topsep=-10000pt]}{\end{itemize}}

\setlength{\headheight}{20.1pt}
\setlength{\headsep}{10pt}

\renewcommand{\familydefault}{\sfdefault} % everything sans-serif

% font size -> set to equivalent to footnotesize of 10 pt
% it affect only the \\normalsize command, all are still relative to default's
% article one (10pt)
\newcommand\articlenormalsize{\fontsize{10pt}{12pt}\selectfont}
\renewcommand\normalsize{\footnotesize}

% redefine sections padding
\makeatletter
\renewcommand{\section}{\@startsection{section}{1}{0mm}%
                                {-1ex plus -.5ex minus -.2ex}%
                                {1ex}%
                                {\normalfont\articlenormalsize\bfseries}}
\renewcommand{\subsection}{\@startsection{subsection}{2}{4mm}%
                                {-1ex}%
                                {0.5ex}%
                                {\normalfont\articlenormalsize\bfseries}}
\renewcommand{\subsubsection}{\@startsection{subsubsection}{3}{16mm}%
                                {-1ex plus -.5ex minus -.2ex}%
                                {1ex plus .2ex}%
                                {\normalfont\small\bfseries}}
\makeatother

%\pagestyle{empty}
% set section counter to 0
\setcounter{secnumdepth}{0}

% Text at right
%\raggedright

\renewcommand{\baselinestretch}{0.9}

% Set up the interline space in tables
\renewcommand{\arraystretch}{0.0}
\linespread{0.85}

% https://github.com/gpoore/minted/issues/82
\makeatletter
\newenvironment{catala}{%
  \VerbatimEnvironment
  \let\FV@ListVSpace\relax
  \begin{minted}[autogobble]{catala_fr}}%
 {\end{minted}}
\makeatother

\setlength{\columnsep}{12pt}
\setlength{\columnseprule}{.1pt}

\begin{document}

\makeheader{\MakeUppercase{La syntaxe de Catala} \hspace{2em}\small version française}{v0.10.0}{1}
% warning: this is repeated for p.2
% \begin{center}
%   2022-12-15
%   \hfill
%   {\Huge\bfseries\sffamily
%     Catala Syntax Cheat Sheet}
%   \hfill
%   v.0.8.0
% \end{center}
% \hrule
\null
\begin{multicols}{3}

\newlength\cola
\setlength\cola{\dimexpr .65\columnwidth -2\tabcolsep}
\newlength\colb
\setlength\colb{\dimexpr .35\columnwidth -2\tabcolsep}

\section{Programmation littéraire}

\begin{tabular}{@{}p{\cola}>{\slshape}p{\colb}@{}}
  \begin{catala}
    # Titre
    ### Sous-sous-titre
  \end{catala}
  & En-têtes
  \\
  \begin{catala}
    # Article 1 | JORFARTI000012345678
    # Article 2 | LEGIARTI000012345678
    # Décision 3 | CETATEXT000012345678
  \end{catala}
  & Référence au journal officiel
  \\
  \begin{catala}
    ```catala      ```catala-metadata
    ```            ```
  \end{catala}
  & Bloc de code / métadonnées
  \\
  \begin{catala}
    > Module Mdl
  \end{catala}
  & Déclaration de mod.
  \\
  \begin{catala}
    > Usage de Mdl en tant que M
  \end{catala}
  & Import de module
  \\
  \begin{catala}
    > Inclusion: foo.catala_en
  \end{catala}
  & Inclusion textuelle
  \\
\end{tabular}

\newcommand*\FancyVerbStartString{```catala}
\newcommand*\FancyVerbStopString{```}

\section{Littéraux et types}

\begin{tabular}{@{}p{\cola}p{\colb}@{}}

  \begin{catala}
    ```catala
    vrai                faux
    ```
  \end{catala}
  & \begin{catala}
    ```catala
    booléen
    ```
  \end{catala}
  \\
  \begin{catala}
    ```catala
    65536
    ```
  \end{catala}
  & \begin{catala}
    ```catala
    entier
    ```
  \end{catala}
  \\
  \begin{catala}
      ```catala
      65536,262144        37%
      ```
    \end{catala}
  & \begin{catala}
    ```catala
    décimal
    ```
  \end{catala}
  \\
  \begin{catala}
    ```catala
    1 234 567,89€
    ```
  \end{catala}
  & \begin{catala}
    ```catala
    argent
    ```
  \end{catala}
  \\
  \begin{catala}
    ```catala
    |2024-04-01|
    ```
  \end{catala}
  & \begin{catala}
    ```catala
    date
    ```
  \end{catala}
  \\
  \begin{catala}
    ```catala
    254 jour     4 mois      1 an
    ```
  \end{catala}
  & \begin{catala}
    ```catala
    durée
    ```
  \end{catala}
  \\
  \begin{catala}
    ```catala
    [ 12; 24; 36 ]
    ```
  \end{catala}
  & \begin{catala}
    ```catala
    liste de entier
    ```
  \end{catala}
  \\
  \begin{catala}
    ```catala
    (|2024-04-01|, 30€, 1%)
    ```
  \end{catala}
  & \begin{catala}
    ```catala
    (date,argent,décimal)
    ```
  \end{catala}
  \\
  \begin{catala}
    ```catala
    f de x, y égal à
      x * y / 12,0€
    ```
  \end{catala}
  & \begin{catala}
    ```catala
    décimal dépend de
      x contenu argent,
      y contenu décimal
    ```
  \end{catala}

  \\

  \begin{catala}
    ```catala
    Struct1 { -- chp1: 9 -- chp2: 7% }
    ```
  \end{catala}
  & \begin{catala}
    ```catala
    Struct1
    ```
  \end{catala}

  \\

  \begin{catala}
    ```catala
    Cas1 contenu 12        Cas2
    ```
  \end{catala}
  & \begin{catala}
    ```catala
    Énum1
    ```
  \end{catala}
  \\
\end{tabular}

\section{Opérations}

\begin{tabular}{@{}p{\cola}>{\slshape}p{\colb}@{}}
  \begin{catala}
    ```catala
    non a         a et b
    a ou b        # "ou à défaut"
    a ou bien b   # ou exclusif
    ```
  \end{catala}
  & Opérateurs logiques
  \\
  \begin{catala}
    ```catala
    - a          a + b        a - b
    a * b        a / b
    ```
  \end{catala}
  & Arithmétique
  \\
  \begin{catala}
    ```catala
    a = b        a < b        a <= b
    a != b       a > b        a >= b
    ```
  \end{catala}
  & Comparaisons
  \\
  \begin{catala}
    ```catala
    décimal de 44
    argent de 23,15
    ```
  \end{catala}
  & Conversions
  \\
  \begin{catala}
    ```catala
    arrondi de 9,99€
    ```
  \end{catala}
  & Arrondis
  \\
  \begin{catala}
    ```catala
    accès_année de ...
    premier_jour_du_mois de ...
    ```
  \end{catala}
  & Éléments de dates
  \\
  \begin{catala}
    ```catala
    a +! b   a +. b   a +€ b   a +^ b
    # entier décimal  argent   durée
    ```
  \end{catala}
  & Opérateurs à types explicites
  \\
\end{tabular}

\section{Déclaration des métadonnées}

\begin{tabular}{@{}p{\cola}>{\slshape}p{\colb}@{}}


  \begin{catala}
    ```catala
    déclaration structure Struct1:
      donnée chp1 contenu entier
      donnée chp2 contenu décimal
    ```
  \end{catala}
  & Déclaration de structure
  \\
  \begin{catala}
    ```catala
    déclaration énumération Énum1:
      -- Cas1 contenu entier
      -- Cas2
    ```
  \end{catala}
  & Déclaration d'énumération
  \\
  \begin{catala}
    ```catala
    déclaration champ d'application Chp1:
      interne var1 contenu entier
      interne var2 condition
      sub1 champ d'application Chp0
    ```
  \end{catala}
  & Déclaration de champ d'application
  \\
  \begin{catala}
    ```catala
      interne var1 contenu ...
      résultat var3 contenu ...
      entrée var4 contenu ...
      entrée résultat var5 contenu ...
      contexte var6 contenu ...
      contexte résultat var7 contenu ...
      résultat sub2
        champ d'application Scope0
    ```
  \end{catala}
  & Qualificateurs d'entrée-sortie
  \\
  \begin{catala}
    ```catala
      interne var1 contenu ...
        état avant
        état après
    ```
  \end{catala}
  & Transitions d'état
  \\
  \begin{catala}
    ```catala
    déclaration const contenu décimal
      égal à 17,1
    ```
  \end{catala}
  & Définition globale
  \\
  \begin{catala}
    ```catala
    déclaration carré contenu décimal
      dépend de x contenu décimal
      égal à x * x
    ```
  \end{catala}
  & Définition de fonction globale
  \\
\end{tabular}

\section{Expressions}

\begin{tabular}{@{}p{\cola}>{\slshape}p{\colb}@{}}
  \begin{catala}
    ```catala
    soit x égal à 36 - 5 dans ...
    ```
  \end{catala}
  & Définition locale
  \\
  \begin{catala}
    ```catala
    selon expr sous forme
    -- Cas1 de x : ...
    -- Cas2 : ...
    -- n'importe quel : ...
    ```
  \end{catala}
  & Filtrage par motif
  \\
  \begin{catala}
    ```catala
    expr sous forme Cas1
    expr sous forme Cas1 de x
         et x >= 2
    ```
  \end{catala}
  & Test de filtrage\newline avec variable optionnelle
  \\
  \begin{catala}
    ```catala
    struc1 mais en remplaçant
      { -- fld2: 8% }
    ```
  \end{catala}
  & Remplacement de champs
  \\
  \begin{catala}
    ```catala
    struc1.chp2         tuple1.2
    sub1.var0
    ```
  \end{catala}
  & Champ, élément de n-uplet, var. de s/s-champ
  \\
  \begin{catala}
    ```catala
    f de 44,50€, 1/3
    ```
  \end{catala}
  & Appel de fonction
  \\
  \begin{catala}
    ```catala
    résultat de Chp1
      avec { -- chp1: 9 -- chp2: 15% }
    ```
  \end{catala}
  & Appel direct de champ d'application
  \\
  \begin{catala}
    ```catala
    si ... alors ... sinon ...
    ```
  \end{catala}
  & Branchement
  \\
  \begin{catala}
    ```catala
    var1 état avant
    ```
  \end{catala}
  & Accès à un état
  \\
\end{tabular}

\columnbreak

\section{Définition de champ d'application}

\begin{tabular}{@{}p{\cola}>{\slshape}p{\colb}@{}}
  \begin{catala}
    ```catala
    champ d'application Chp1: ...
    ```
  \end{catala}
  & Utilisation
  \\
  \begin{catala}
    ```catala
    champ d'application Chp1
      sous condition var1 >= 2: ...
    ```
  \end{catala}
  & Avec condition générale
  \\
  \begin{catala}
    ```catala
    définition var1 égal à ...
    ```
  \end{catala}
  & Déf. inconditionnelle
  \\
  \begin{catala}
    ```catala
    définition var1
      sous condition ...
      conséquence égal à ...
    ```
  \end{catala}
  & Définition conditionnelle
  \\
  \begin{catala}
    ```catala
    règle var2
      sous condition var1 >= 2
      conséquence rempli
    ```
  \end{catala}
  & Règle\newline(définition de condition)
  \\
  \begin{catala}
    ```catala
      conséquence non rempli
    ```
  \end{catala}
  & Règle négative
  \\
  \begin{catala}
    ```catala
    définition f de x, y égal à ...
    ```
  \end{catala}
  & Déf./règle fonction
  \\
  \begin{catala}
    ```catala
    étiquette étq1 définition var1 ...
    ```
  \end{catala}
  & Déf./règle étiquetée
  \\
  \begin{catala}
    ```catala
    exception étq1 définition var1 ...
    ```
  \end{catala}
  & Exc. à déf. étiquetée
  \\
  \begin{catala}
    ```catala
    exception définition var1 ...
    ```
  \end{catala}
  & Exception à implicite
  \\
  \begin{catala}
    ```catala
    définition var1
      état avant
      égal à ...
    ```
  \end{catala}
  & Définition d'états
  \\
  \begin{catala}
    ```catala
    assertion ...
    ```
  \end{catala}
  & Assertion
  \\
  \begin{catala}
    ```catala
    date arrondi décroissant
    ```
  \end{catala}
  & Mode arrondi dates
  \\
\end{tabular}

\vspace{-5pt}

\vfill

\section{Opérations sur les listes}

\vspace{-5pt}

\begin{tabular}{@{}p{\cola}>{\slshape}p{\colb}@{}}
  \begin{catala}
    ```catala
    lst contient 3
    ```
  \end{catala}
  & Test de présence
  \\
  \begin{catala}
    ```catala
    existe x parmi lst tel que x > 2
    ```
  \end{catala}
  & Test d'existence
  \\
  \begin{catala}
    ```catala
    pour tout x parmi lst on a x > 2
    ```
  \end{catala}
  & Test pour tout
  \\
  \begin{catala}
    ```catala
    transforme chaque x parmi lst
      en x - 2
    ```
  \end{catala}
  & Application un-à-un
  \\
  \begin{catala}
    ```catala
    liste de x parmi lst tel que x > 2
    ```
  \end{catala}
  & Filtrage
  \\
  \begin{catala}
    ```catala
    transforme chaque x parmi lst
      tel que x > 2 en x - 2
    ```
  \end{catala}
  & Filtrage + application
  \\
  \begin{catala}
    ```catala
    transforme chaque (x, y)
      parmi (lst1, lst2) en x + y
    ```
  \end{catala}
  & Application multiple
  \\
  \begin{catala}
    ```catala
    lst1 ++ lst2
    ```
  \end{catala}
  & Réunion
  \\
  \begin{catala}
    ```catala
    somme entier de lst
    ```
  \end{catala}
  & Agrégation
  \\
  \begin{catala}
    ```catala
    nombre de lst
    ```
  \end{catala}
  & Comptage
  \\
  \begin{catala}
    ```catala
    maximum de lst
      ou si liste vide alors -1
    ```
  \end{catala}
  & Extremums\newline (optionnel: défaut)
  \\
  \begin{catala}
    ```catala
    contenu de x parmi lst
      tel que x * x est minimum
      ou si liste vide alors -1
    ```
  \end{catala}
  & Élément selon extremum\newline (optionnel: défaut)
  \\
  \begin{catala}
    ```catala
    combine tout x parmi lst
      dans acc initialement 0
      avec acc + x
    ```
  \end{catala}
  & Accumulation
  \\
\end{tabular}

\end{multicols}
\end{document}
